\documentclass[../main.tex]{subfiles}
\graphicspath{{\currfiledir}}
\begin{document}
\chapter{Terminology}

\section{Grapheme}
\textcquote[p.204]{crystal2010}{Graphemes are the smallest units in a writing system capable of 
causing a contrast in meaning.}
In linguistics, graphemes are often placed into angle brackets, e.g. \grapheme{a} or \grapheme{b}.
Sometimes \emph{graphemes} are called \emph{signs}.
The term \emph{sign} and \emph{grapheme} are considered to be equivalent and 
interchangeable throughout this work.

\section{Allograph}
\textcquote[p.204]{crystal2010}{Graphemes are abstract units which may adopt a variety of forms 
\elide Each of these possible forms is known as \emph{graphs}\elide
There is a vast amount of physical variation in the shapes of graphs that does not affect the 
underlying identity of the grapheme\elide 
When graphs are analyzed as variants of a grapheme, they are known as \emph{allographs}.}
The Maya script uses a lot of allographs.
For example, the syllable \syllable{u} can be written in many ways all having the same meaning.
See \ref{fig:terminology-grapheme-u-allographs} for a selection of allographs.
\begin{center}
    \includegraphics[width=\textwidth,keepaspectratio]{img/grapheme-u-allographs}
    \captionof{figure}{Some allographs of the grapheme \grapheme{u}}
    \label{fig:terminology-grapheme-u-allographs}
\end{center}

Tikal Temple IV, Lintel 2 A7


\end{document}
